\heading{Homework 1 · Part A}{Formative mock examination · 120 minutes · 100 points}
Work on this part honestly in one 120-minute sitting. Write by hand on paper. Use only this question paper; no notes, books, internet or AI during the diagnostic sitting. This is a recommended practice protocol, not an announcement of the actual final exam's permitted materials. After the timed attempt, you may study the material and clearly label any later corrections. Preserve the original attempt for feedback.

The TA will mark this part seriously so that you can assess readiness for the exam. Mistakes in Part A do \textbf{not} reduce your final course grade. Its points are diagnostic and are not added to the assessed Part B score. This mock covers Lecture 1, not the whole final-exam syllabus. Show reasoning: an unsupported true/false label earns no credit.
\conventions
\task{A1. True or false, with justification}{40 points · suggested 40 minutes}
Each statement is worth 5 points. Unless another domain is specified, use $X=\R_+^2$. State whether each claim is true or false. Prove a true claim and give a fully specified counterexample for a false claim.
\begin{enumerate}[label=(\alph*)]
\item The relation $x\pref y\iff x_1+x_2\ge y_1+y_2-1$ is complete.
\item The relation $x\pref y\iff x_1+x_2\ge y_1+y_2-1$ is transitive.
\item The preferences represented by $u(x)=x_1+\mathbf1\{x_2\ge1\}$ are continuous.
\item The preferences represented by $u(x)=\min\{x_1,2x_2\}$ are convex.
\item Continuous preferences imply that \emph{every} representing utility function is continuous.
\item A continuous utility function must represent convex preferences.
\item If $u$ is concave on a convex domain, the preferences it represents are convex.
\item Rational, continuous, strictly convex preferences must be locally nonsatiated.
\end{enumerate}
\clearpage
\task{A2. A consumer with an ideal bundle}{30 points · suggested 35 minutes}
Let $X=\R_+^2$ and $u(x)=-(x_1-3)^2-(x_2-2)^2$.
\begin{enumerate}[label=(\alph*)]
\item Draw the indifference set through $(3,0)$ and the weak upper contour set of $(3,0)$. Respect the domain. \hfill [6]
\item Prove rationality and continuity, and determine whether the preferences are convex or strictly convex. A picture alone is not a proof. \hfill [8]
\item Are these preferences monotone, strongly monotone or LNS? Give a witness for each failed property. \hfill [6]
\item Let $v=e^u$ and $w=\max\{u,-1\}$. Which represents the same preferences as $u$? If not, exhibit two bundles whose comparison changes. Does the ranking represented by $v$ remain strictly convex? \hfill [6]
\item Could \emph{any} convex utility function represent these preferences on $X$? Use two carefully chosen indifferent bundles. \hfill [4]
\end{enumerate}
\clearpage
\task{A3. A majority of criteria}{15 points · suggested 20 minutes}
On $X=\R^3$, define $x\pref y$ if at least two of the three inequalities $x_i\ge y_i$ hold.
\begin{enumerate}[label=(\alph*)]
\item Prove completeness, including coordinate ties. \hfill [4]
\item Test transitivity using $a=(4,2,0)$, $b=(2,0,4)$ and $c=(0,4,2)$. State the required comparisons explicitly. \hfill [5]
\item Is the relation continuous under the definition on this paper? \hfill [3]
\item Could any real-valued utility represent it? Which theorem or elementary implication decides this? \hfill [3]
\end{enumerate}
\task{A4. A qualitative question}{15 points · suggested 20 minutes}
A researcher writes: ``Our consumer has complete, transitive preferences. Therefore we can write a real-valued utility function. After an increasing transformation, we may assume it is continuous and concave. Thus every choice problem has a unique solution.''

Assess this argument step by step. Identify which conclusions need additional assumptions, give at least two explicit counterexamples, and state a correct set of sufficient assumptions for existence and uniqueness on a nonempty compact convex feasible set in $\R_+^2$. Explain why continuity and concavity are different issues. Do not simply list theorem names.

\medskip Total suggested time: 115 minutes; keep 5 minutes to check your work.
